% Source: https://artofproblemsolving.com/community/c3h2615296_nice_problem
% Stable identity: aops-topic-2615296-q01
Let $n=b-a+1\ge1$ be the number of terms and $m=a+b$. The sum is $\dfrac{n(a+b)}{2}$, so
\[nm=2\cdot10!=2^9\cdot3^4\cdot5^2\cdot7.\]
Here $m>0$, and $m-n=2a-1$ is odd, so $n$ and $m$ have opposite parity. Conversely, if $nm=2\cdot10!$ with $n,m$ positive integers of opposite parity, then $a=\dfrac{m-n+1}{2}$ and $b=\dfrac{m+n-1}{2}$ are integers with $a\le b$, and the sum from $a$ to $b$ is $10!$.

Since $nm$ contains $2^9$ and exactly one of $n$, $m$ is even, the even one contains all nine factors of $2$, and the odd one is an odd divisor of $10!$. There are $(4+1)(2+1)(1+1)=30$ odd divisors of $10!=2^8\cdot3^4\cdot5^2\cdot7$, and either $n$ or $m$ can be the odd one. This gives $2\cdot30=\boxed{60}$ pairs.

(For example, $n=1$ gives the single term $a=b=10!$, and the ``mirror'' pair with $n=2\cdot10!$ is $a=1-10!$, $b=10!$.)
